\section{Detailed Result Ledger}
\label{app:results}

This appendix records the source-backed result tables used by the main paper.
All paths are relative to the repository root.

\subsection{Phase 1: Cross-Domain Structure Compatibility}

Source:
\path{experiments/structure_compatible_generalization/results/structure_compatible_l4_2026_07_06.md}.

\begin{table*}[h]
\centering
\small
\begin{tabular}{lrrrr}
\toprule
Domain & rows & mean train & mean ID & mean OOD \\
\midrule
modular exact & 2 & 1.000 & 1.000 & 0.545 \\
modular neural & 128 & 0.732 & 0.732 & 0.165 \\
symbolic cyclic & 128 & 0.973 & 0.984 & 0.394 \\
vision rotation & 96 & 0.773 & 0.773 & 0.380 \\
\bottomrule
\end{tabular}
\caption{Phase 1 domain sizes and mean accuracies.}
\end{table*}

\begin{table*}[h]
\centering
\small
\begin{tabular}{lrrrr}
\toprule
Selector & domains & mean selected OOD & note \\
\midrule
true compatibility & 3 & 1.000 & oracle structure \\
ID validation accuracy & 3 & 0.660 & standard selector \\
train accuracy & 3 & 0.660 & standard selector \\
inferred compatibility & 2 & 0.531 & unavailable in one domain \\
wrong compatibility & 3 & 0.458 & negative control \\
parameter L2 & 3 & 0.426 & proxy baseline \\
sharpness proxy & 3 & 0.410 & proxy baseline \\
negative train loss & 3 & 0.117 & proxy baseline \\
\bottomrule
\end{tabular}
\caption{Phase 1 OOD-free selection by predictor family.}
\end{table*}

\subsection{Phase 2: Inferred Transformations and Intervention}

Source:
\path{experiments/structure_compatible_generalization/results/phase2_transformations_2026_07_06.md}.

The main regime transition is from post-hoc oracle compatibility to supported
modular shifts inferred from observed train-label overlaps. The evidence-bearing
result is regularization, not the one-domain selector sanity check.

\begin{table*}[h]
\centering
\small
\begin{tabular}{rrrrrr}
\toprule
reg. & N & mean train & mean OOD & high-ID N & high-ID OOD \\
\midrule
0.000 & 180 & 0.719 & 0.134 & 98 & 0.178 \\
0.050 & 180 & 0.719 & 0.320 & 98 & 0.519 \\
0.200 & 180 & 0.723 & 0.352 & 99 & 0.573 \\
\bottomrule
\end{tabular}
\caption{Phase 2 compatibility regularization improves OOD while preserving a
high-ID comparison surface.}
\end{table*}

\subsection{Phase 3: Learned Generators and Vision Transfer}

Source:
\path{experiments/structure_compatible_generalization/results/phase3_learned_generators_2026_07_06.md}.

\begin{table*}[h]
\centering
\small
\begin{tabular}{lrrrr}
\toprule
Domain & predictor & Pearson $r$ & Spearman $r$ & N \\
\midrule
modular learned generator & true compatibility & 0.812 & 0.367 & 270 \\
modular learned generator & discovered compatibility & 0.787 & 0.474 & 270 \\
modular learned generator & ID validation & 0.398 & 0.120 & 270 \\
modular learned generator & wrong compatibility & -0.085 & 0.049 & 270 \\
vision learned generator & true compatibility & 0.678 & 0.591 & 144 \\
vision learned generator & discovered compatibility & 0.599 & 0.502 & 144 \\
vision learned generator & ID validation & 0.661 & 0.632 & 144 \\
\bottomrule
\end{tabular}
\caption{Phase 3 predictor correlations. The vision arm is a mixed case where
ID remains strong, so the paper should not overstate compatibility dominance.}
\end{table*}

\begin{table*}[h]
\centering
\small
\begin{tabular}{lrrr}
\toprule
Vision regime & N & mean OOD & mean learned compatibility \\
\midrule
none & 36 & 0.258 & 0.486 \\
learned augmentation & 36 & 0.644 & 0.793 \\
oracle augmentation & 36 & 0.693 & 0.801 \\
random augmentation & 36 & 0.610 & 0.760 \\
\bottomrule
\end{tabular}
\caption{Phase 3 vision transfer arm. The random augmentation result means the
vision claim should be framed as transfer evidence with a nontrivial control,
not as a pure learned-generator victory.}
\end{table*}

\subsection{Phase 4: Language-Like Template Substitution}

Source:
\path{experiments/structure_compatible_generalization/results/language_template_substitution_2026_07_06.md}.

\begin{table*}[h]
\centering
\small
\begin{tabular}{lrrr}
\toprule
Predictor & Pearson $r$ & Spearman $r$ & N \\
\midrule
discovered compatibility & 0.858 & 0.596 & 256 \\
true compatibility & 0.853 & 0.628 & 256 \\
negative train loss & 0.492 & 0.232 & 256 \\
ID validation & 0.478 & 0.127 & 256 \\
train accuracy & 0.478 & 0.127 & 256 \\
wrong compatibility & -0.088 & 0.059 & 256 \\
sharpness proxy & -0.111 & 0.063 & 256 \\
parameter L2 & -0.476 & -0.189 & 256 \\
\bottomrule
\end{tabular}
\caption{Phase 4 finite-text bridge.}
\end{table*}

\begin{table*}[h]
\centering
\small
\begin{tabular}{rrrrrrr}
\toprule
reg. & N & train & ID & OOD & high-ID N & high-ID OOD \\
\midrule
0.000 & 64 & 0.554 & 0.554 & 0.147 & 29 & 0.224 \\
0.050 & 64 & 0.553 & 0.553 & 0.234 & 29 & 0.415 \\
0.200 & 64 & 0.555 & 0.555 & 0.255 & 29 & 0.463 \\
0.500 & 64 & 0.558 & 0.558 & 0.277 & 29 & 0.512 \\
\bottomrule
\end{tabular}
\caption{Phase 4 compatibility regularization.}
\end{table*}

\subsection{Phase 5: Semantic Retrieval Transfer}

Source:
\path{experiments/structure_compatible_generalization/results/semantic_retrieval_transfer_2026_07_06.md}.

\begin{table*}[h]
\centering
\small
\begin{tabular}{lrrr}
\toprule
Predictor & Pearson $r$ & Spearman $r$ & N \\
\midrule
true compatibility & 0.940 & 0.915 & 64 \\
discovered compatibility & 0.861 & 0.863 & 64 \\
train accuracy & 0.772 & 0.628 & 64 \\
ID validation & 0.709 & 0.488 & 64 \\
wrong compatibility & -0.872 & -0.919 & 64 \\
\bottomrule
\end{tabular}
\caption{Phase 5 frozen-encoder semantic retrieval.}
\end{table*}

\subsection{Phase 6: Semantic Selection Control}

Source:
\path{experiments/structure_compatible_generalization/results/semantic_selection_control_2026_07_06.md}.
Bootstrap source:
\path{experiments/structure_compatible_generalization/results/semantic_selection_bootstrap_2026_07_06.md}.
Row release:
\path{experiments/structure_compatible_generalization/results/semantic_selection_row_release_2026_07_06.md}.
Tie-break stress:
\path{experiments/structure_compatible_generalization/results/semantic_selection_tiebreak_stress_2026_07_06.md}.

All registered selector gates passed: minimum zoo count, beats ID validation,
beats train accuracy, beats random candidate, wrong control fails, and accepted.
The table in the main paper reports the full selector comparison. The
regenerated 120-zoo payload gives learned-minus-random selected OOD 0.059 with
95\% CI [0.052, 0.065], learned-minus-ID 0.059 [0.052, 0.065],
learned-minus-wrong 0.227 [0.221, 0.233], and learned regret 0.014
[0.011, 0.018]. The row release tracks 1,440 diagnostic rows and 840 selection
records. Mean-tie, worst-tied-candidate, and random-tie bootstrap variants all
pass the registered selector gates.

\section{Companion Benchmark Appendix}
\label{app:agents}

The companion benchmark is not the main ICML claim, but it is useful because it
shows the same anti-proxy rule in finite agents.

\subsection{Minimum Pass Rule}

Source: \path{docs/causally_grounded_agents_benchmark.md}.

A suite does not pass on behavior alone. A pass requires:
\begin{enumerate}
  \item the task behavior gate,
  \item the relevant causal-structure gate, and
  \item anti-cheat controls specific to that suite.
\end{enumerate}

The score should remain vector-valued: behavior, causal representation,
attribution, inquiry, commitment, and generalization.

\subsection{Suite C: Re-Engagement Under World Change}

Sources:
\path{experiments/world_responds/results/suite_c_reengagement_2026_07_06.md}
and
\path{experiments/world_responds/results/suite_c_neural_transfer_2026_07_06.md}.

The hand-policy headline condition is \texttt{burst\_then\_refractory}. It
passes C1-C6 with final affected MAE 0.112, first-shift selectivity 16.927,
second-shift reopenability 16.667, and 22.6 probes. The signal-layer
\texttt{fixed\_surprise\_decrement} control fails correctly: it looks quiet but
ends at final MAE 0.517.

The learned probe head passes the same held-out C1-C6 gate with final affected
MAE 0.112, first-shift selectivity 16.667, second-shift reopenability 17.448,
and 23.1 probes. Stale-signal, wrong-signal, signal-suppression, and
matched-random controls fail in the intended way.

\subsection{Suite D/E: Long-Horizon Commitment Surfaces}

Source: \path{experiments/long_horizon_bottleneck/BENCHMARK_CARD.md}.

The long-horizon benchmark asks whether a future-critical variable survives the
surfaces where it must control action: hidden state, generated JSON, tool
arguments, repair branches, and value-token causal readouts. The strongest
lesson is the commitment-surface law: memory becomes agent-relevant when it is
bound to a future action or tool surface.

\section{Reproducibility Notes}

The source reports list base seeds, model/config counts, GPU type, and budget
checks. The paper package should be submitted with:
\begin{itemize}
  \item the six source result reports above,
  \item figure regeneration scripts under
  \path{scripts/build_structure_compatible_*.py} and
  \path{scripts/export_structure_compatible_artifacts.py},
  \item task code under
  \path{experiments/structure_compatible_generalization/},
  \item public benchmark schema under
  \path{docs/causally_grounded_agents_release_schema.json}.
\end{itemize}

\section{Reviewer-Facing Weaknesses}

The strongest reviewer objections are:
\begin{enumerate}
  \item Some phases use oracle or partly hand-designed transformation families.
  \item The vision arm has a strong ID baseline and a strong random
  augmentation control.
  \item The semantic tasks are constructed around frozen encoders and model
  zoos; they are not broad natural-language deployment.
  \item Phase 6 now has bootstrap intervals and tie-break stress tests, but
  earlier phases have only smoke row fixtures unless full Modal payloads are
  restored.
  \item The finite-agent benchmark material should remain an appendix unless a
  separate benchmark paper is submitted.
\end{enumerate}

\section{Next Experiments}

The next ICML-strengthening experiments are:
\begin{itemize}
  \item paired-difference presentation for the existing Phase 6 selector rows,
  \item larger semantic model zoos with a third frozen encoder,
  \item held-out transformation-family discovery where candidate transforms are
  learned from one family and evaluated on another,
  \item an ablation that holds augmentation volume fixed while varying
  transformation correctness,
  \item less privileged source-estimation and open-agent transfer for the
  teacher-free Suite C companion benchmark.
\end{itemize}
